\documentclass[aspectratio=169,10pt,fontset=fandol]{ctexbeamer}
\usepackage{amsmath,amssymb,mathtools,bm,mathrsfs}
\usepackage{graphicx,booktabs,array,tabularx,multirow}
\usepackage{tikz,pgfplots,xparse,xstring,mdframed}
\usetikzlibrary{arrows.meta,calc,angles,positioning,shapes.geometric,graphs}
\pgfplotsset{compat=1.18}

\setbeamerfont{normal text}{size=\normalsize}
\setbeamerfont{title}{size=\large}
\setbeamerfont{subtitle}{size=\normalsize}
\setbeamerfont{frametitle}{size=\large}
\setbeamerfont{framesubtitle}{size=\small}
\setbeamerfont{block title}{size=\normalsize}
\setbeamerfont{block body}{size=\normalsize}
\setbeamerfont{section in toc}{size=\normalsize}
\setbeamerfont{section in head/foot}{size=\footnotesize}
\setbeamerfont{author in head/foot}{size=\footnotesize}
\setbeamerfont{title in head/foot}{size=\footnotesize}
\setbeamerfont{date in head/foot}{size=\footnotesize}
\setbeamertemplate{navigation symbols}{}
\setbeamertemplate{itemize item}{\normalsize\textbullet}
\setbeamertemplate{itemize subitem}{\small\textbullet}
\setbeamercolor{structure}{fg=blue!65!black}
\setbeamercolor{block title}{fg=blue!65!black,bg=blue!8}
\setbeamercolor{block body}{fg=black,bg=blue!3}
\setbeamertemplate{footline}{\leavevmode\hbox to\paperwidth{\hfill{\footnotesize\insertframenumber}\hspace*{1em}}\vskip3pt}

% ctex fontset=fandol; system fonts are not redistributed.

\newenvironment{proofdetail}{\par\noindent\textbf{证明：}\quad}{\par}
\renewenvironment{solution}{\par\noindent\textbf{解：}\quad}{\par}
\makeatletter
\@ifundefined{definition}{\newtheorem{definition}{定义}[section]}{}
\@ifundefined{theorem}{\@ifundefined{c@definition}{\newtheorem{theorem}[theorem]{定理}}{\newtheorem{theorem}[definition]{定理}}}{}
\@ifundefined{exercise}{\theoremstyle{definition}\@ifundefined{c@definition}{\newtheorem{exercise}[theorem]{练习}}{\newtheorem{exercise}[definition]{练习}}\theoremstyle{plain}}{}
\@ifundefined{exercise}{\theoremstyle{definition}\newtheorem{exercise}[theorem]{练习}\theoremstyle{plain}}{}
\makeatother
\makeatletter
\newcommand{\mathagentsectionlabelcontext}{%
\protected@edef\@currentlabel{\thesection}%
\edef\@currentHref{Outline\the\c@part.\the\c@section}%
}
\makeatother
\title{讲义兼容}
\begin{document}
\section{第一节}
% FRAME-ID F001
\typeout{MATH-AGENT-FRAME-BEGIN F001}
\begin{frame}[t,fragile]{目录}
\normalsize
{\color{blue}\normalsize 1. 第一节\par}\vspace{0.22em}
{\color{gray!40}\normalsize 2. 第二节\par}\vspace{0.22em}
{\color{gray!40}\normalsize 3. 练习答案\par}\vspace{0.22em}
\end{frame}
\typeout{MATH-AGENT-FRAME-END F001}
% FRAME-ID F002
\typeout{MATH-AGENT-FRAME-BEGIN F002}
\begin{frame}[t,fragile]{第一节}
\normalsize
{\normalsize
% PRESENTATION-SCAFFOLD original heading label binds to section
\mathagentsectionlabelcontext
% SOURCE-BEGIN U0002 0 49

\label{sec:first}
第一节内容；第二节见第\ref{sec:second}节。

% SOURCE-END U0002
}
\end{frame}
\typeout{MATH-AGENT-FRAME-END F002}
% FRAME-ID F003
\typeout{MATH-AGENT-FRAME-BEGIN F003}
\begin{frame}[t,fragile]{第一节}
\normalsize
{\normalsize
% SOURCE-BEGIN U0003 0 59
\begin{exercise}[求和]\label{ex:first}计算 $1+1$。\end{exercise}
% SOURCE-END U0003
}
\end{frame}
\typeout{MATH-AGENT-FRAME-END F003}
\section{第二节}
% FRAME-ID F004
\typeout{MATH-AGENT-FRAME-BEGIN F004}
\begin{frame}[t,fragile]{目录}
\normalsize
{\color{black}\normalsize 1. 第一节\par}\vspace{0.22em}
{\color{blue}\normalsize 2. 第二节\par}\vspace{0.22em}
{\color{gray!40}\normalsize 3. 练习答案\par}\vspace{0.22em}
\end{frame}
\typeout{MATH-AGENT-FRAME-END F004}
% FRAME-ID F005
\typeout{MATH-AGENT-FRAME-BEGIN F005}
\begin{frame}[t,fragile]{第二节}
\normalsize
{\normalsize
% PRESENTATION-SCAFFOLD original heading label binds to section
\mathagentsectionlabelcontext
% SOURCE-BEGIN U0004 0 49

\label{sec:second}
第二节内容；第一节见第\ref{sec:first}节。

% SOURCE-END U0004
}
\end{frame}
\typeout{MATH-AGENT-FRAME-END F005}
% FRAME-ID F006
\typeout{MATH-AGENT-FRAME-BEGIN F006}
\begin{frame}[t,fragile]{第二节}
\normalsize
{\normalsize
% SOURCE-BEGIN U0005 0 56
\begin{exercise}\label{ex:second}计算 $2+2$。\end{exercise}
% SOURCE-END U0005
}
\end{frame}
\typeout{MATH-AGENT-FRAME-END F006}
\section{练习答案}
% FRAME-ID F007
\typeout{MATH-AGENT-FRAME-BEGIN F007}
\begin{frame}[t,fragile]{目录}
\normalsize
{\color{black}\normalsize 1. 第一节\par}\vspace{0.22em}
{\color{black}\normalsize 2. 第二节\par}\vspace{0.22em}
{\color{blue}\normalsize 3. 练习答案\par}\vspace{0.22em}
\end{frame}
\typeout{MATH-AGENT-FRAME-END F007}
% FRAME-ID F008
\typeout{MATH-AGENT-FRAME-BEGIN F008}
\begin{frame}[t,fragile]{练习答案}
\normalsize
{\normalsize
% SOURCE-BEGIN U0006 0 97

\hyperref[ex:first]{练习\ref*{ex:first}的解答：}$2$。
\hyperref[ex:second]{练习\ref*{ex:second}的解答：}$4$。

% SOURCE-END U0006
}
\end{frame}
\typeout{MATH-AGENT-FRAME-END F008}
\end{document}
